\section{Trabajos Relacionados}

La evolución técnica y arquitectónica del Portal Agro-Comercial no ocurre en un vacío académico, sino que dialoga activamente con un cuerpo de conocimiento en ingeniería de software cada vez más sofisticado. En la actualidad, la industria se enfrenta a una paradoja: la necesidad de incrementar la velocidad de entrega (\textit{Delivery Speed}) choca frontalmente con la exigencia de mantener una estabilidad sistémica absoluta (\textit{System Reliability}).

Para fundamentar la propuesta de mejora arquitectónica presentada en este artículo, se trascendió la revisión teórica superficial para realizar un análisis crítico, sistemático y temático de la literatura reciente, cubriendo una ventana de observación de nueve años (2016--2025). Este análisis no se limitó a la enumeración de definiciones, sino que diseccionó los hallazgos de 20 investigaciones clave para entender cómo los gigantes tecnológicos han resuelto problemas que hoy empezamos a ver en el desarrollo de software para el agro. El estado del arte se estructuró bajo cinco lentes temáticos.

\subsection{Estrategias de Despliegue y Configuración Dinámica}

La transición de aplicaciones web monolíticas y estáticas hacia arquitecturas vivas y dinámicas representa una tendencia irreversible en la ingeniería moderna. En este sentido, \cite{meyer2016continuous} son categóricos al establecer, mediante un estudio empírico extensivo, que la adopción de prácticas de entrega continua (\textit{Continuous Delivery}) ha dejado de ser una ventaja competitiva opcional para convertirse en un requisito de supervivencia industrial. Su investigación sugiere que los equipos que despliegan con mayor frecuencia paradójicamente experimentan menos fallos catastróficos. Para un proyecto como el Portal Agro-Comercial, esto implica una refutación del modelo tradicional: el ciclo de ``parar el servidor para actualizar'' no es una medida de seguridad, sino un síntoma de obsolescencia.

Profundizando en esta necesidad de fluidez, \cite{ramaswamy2024zerodowntime} argumentan que, para alcanzar el estándar de oro de ``tiempo de inactividad cero'', es mandatorio desacoplar conceptual y técnicamente el acto de \textit{desplegar} (instalar los binarios en el servidor) del acto de \textit{liberar} (hacer disponible la funcionalidad al usuario). Es en este desacople donde las \textit{Feature Flags} se posicionan no como un truco de programación, sino como el mecanismo arquitectónico clave que habilita el despliegue desacoplado.

Sin embargo, la literatura reciente advierte que esta problemática ha migrado del servidor al cliente. \cite{lee2025frontend} trasladan esta discusión al desarrollo frontend, un aspecto crítico para nuestra plataforma dado el contexto de conectividad intermitente en Teruel. Ellos proponen arquitecturas donde los componentes de la interfaz de usuario (UI) reaccionen dinámicamente a configuraciones remotas, sin necesidad de recargar la página. Para nuestro caso de estudio, esto significaría poder actualizar la visualización de precios o productos críticos sin obligar al agricultor a descargar una nueva versión completa de la aplicación, ahorrando ancho de banda. Complementariamente, \cite{esther2025microservices} refuerzan que, en ecosistemas distribuidos, la inyección de configuraciones dinámicas es lo que garantiza la resiliencia, permitiendo al sistema degradarse elegantemente (desactivando módulos fallidos) en lugar de colapsar ante un error.

\subsection{Complejidad Ciclomática y Costo Cognitivo}

A pesar de sus innegables beneficios operativos, la implementación de \textit{Feature Toggles} conlleva un costo de ingeniería inherente que debe ser gestionado. \cite{rahman2016msr}, en uno de los estudios seminales del campo, advierten sobre el peligro de la proliferación descontrolada de banderas. Su análisis de casos industriales muestra que el uso excesivo de condicionales puede convertir el código base en un laberinto inmanejable.

Esta preocupación teórica ha sido validada empíricamente por investigaciones recientes como la de \cite{sharma2024complexity}, quienes encontraron una correlación directa y positiva entre la densidad de toggles activos y el aumento de la complejidad ciclomática del software. Esto sugiere que cada nuevo toggle incrementa la carga cognitiva necesaria para entender el flujo del programa. Para el Portal Agro-Comercial, la lección es clara: sin una gobernanza adecuada, la solución propuesta podría terminar dificultando el mantenimiento a largo plazo, haciendo el código más difícil de leer para futuros desarrolladores del SENA.

\subsection{El Desafío de la Deuda Técnica: Detección y Limpieza}

De la complejidad no gestionada emerge el concepto de ``Deuda de Toggles'' (\textit{Toggle Debt}), analizado en profundidad por \cite{ortiz2021toggledebt}. El autor insiste en la necesidad de concebir las banderas con un ciclo de vida finito: se debe identificar claramente cuándo una funcionalidad ya es estable y, por tanto, cuándo el toggle debe ser eliminado del código fuente.

Para automatizar este proceso y evitar la acumulación de ``basura lógica'', \cite{li2020capture} y, más recientemente, \cite{wang2025tsdetector}, han propuesto algoritmos de detección estática avanzados. Estas herramientas escanean el Código Fuente Abstracto (AST) para identificar banderas ``zombis'' —aquellas que siempre evalúan a verdadero o falso— en repositorios de código abierto.

La solución industrial a este problema es abordada por \cite{sridharan2020piranha}, quienes presentan la experiencia de Uber con la herramienta Piranha. Su estudio demuestra que la limpieza de código muerto puede ser automatizada refactorizando el código automáticamente. No obstante, el obstáculo principal suele ser humano y cultural; \cite{abdalkareem2021removal} estudiaron las motivaciones psicológicas de los desarrolladores y concluyeron que el miedo a ``romper algo que funciona'' lleva a posponer indefinidamente la eliminación de toggles antiguos, acumulando riesgo silenciosamente hasta que el sistema se vuelve frágil.

\subsection{Interdependencias y Conflictos en Escala}

Al escalar el uso de esta técnica, surgen fenómenos emergentes de segundo orden que no son visibles en proyectos pequeños. \cite{rahman2019chrome} diseccionaron la arquitectura modular de Google Chrome, evidenciando que su capacidad de evolución depende de un sistema masivo de miles de toggles.

Sin embargo, \cite{smith2022interdependencies}, al estudiar el ecosistema de Microsoft Office, resaltan un riesgo crítico: las interdependencias ocultas. Activar un toggle puede requerir implícitamente que otro esté habilitado, creando una red de dependencias que \cite{hal2022interaction} denominan ``interacción de features''. Para el Portal Agro-Comercial, esto implica un reto de diseño: no podemos diseñar las banderas de forma aislada; debemos mapear explícitamente sus relaciones jerárquicas para evitar estados inconsistentes (ej. activar el botón de ``Pago'' sin activar el módulo de ``Cálculo de Impuestos'').

\subsection{Taxonomías y Convergencia con Modelos de Negocio}

Finalmente, para estructurar nuestra propuesta arquitectónica, es vital entender qué tipo de toggle estamos usando. \cite{torres2020differences} realizan una distinción fundamental entre simples opciones de configuración (que son estáticas y de bajo riesgo) y verdaderos \textit{feature flags} (que son dinámicos y alteran el flujo de control).

En la frontera del conocimiento, \cite{dupont2022unifying} proponen unificar los toggles con los modelos de características de las Líneas de Productos de Software (SPL). Esto resulta especialmente relevante para nuestra visión a futuro, donde el portal podría dejar de ser un producto único para convertirse en una familia de productos con variantes para diferentes municipios del Huila.

Asimismo, \cite{chen2025horizon} proponen el marco HORIZON, clasificando los toggles según estrategias de precios. Esta visión es revolucionaria para nuestro contexto, pues sugiere que en el futuro podríamos ofrecer funcionalidades ``Premium'' (como analítica avanzada de cultivos) a grandes cooperativas, gestionando el acceso no mediante versiones diferentes del software, sino mediante banderas activadas por licencia. Finalmente, seguimos las heurísticas de \cite{ajmeri2022heuristics} sobre ``Toggles como Código'' para garantizar una implementación limpia, tipada y mantenible.

\input{tables/toggle_types}
